% Research note: reusable boundary facts; not required for the main theorem.
\paragraph{Grassmann閉包に残る二つの境界平面.}
一般の $G$ に対して、$K_6$ 上の三つの交代形式を同時に零にする三次元部分空間は、
相補的な二つ $U_+,U_-$ である。
これは $G=\operatorname{diag}(I_3,I_3)$ で直接検証できる。
このとき交代形式の値は $v\times v'+w\times w'$ となる。
共通等方的三次元部分空間のいずれかの射影が単射でなければ、
もう一方の射影像が一次元以下となり、その射影の核にある二つの独立なベクトルの外積が非零となって矛盾する。
従って部分空間は $w=Tv$ というグラフであり、
$\bigwedge^2T=-I$、すなわち $T=\pm\sqrt{-1}\,I$ である。
外積二乗写像の微分は可逆なので、この二解は被約である。
射影Grassmann多様体上の零点スキームに対する開性から、同じ性質が一般の $G$ に成立する。

各 $U_\pm$ に対して
\[
 \{U_\pm\oplus\ell:\ell\in\mathbb P(K_3)\}\simeq\mathbb P^2
 \subset X_Q
\]
という線形平面が存在する。
これらの点に対応する行列 $B$ は、行基底変更の後に $\ell(x)e_7$ という行を持つので、
直線 $\ell(x)=0$ の全体で $I_6(L)$ が消える。
従って二つの平面は有限二十一点配置の開部分から除かれる境界である。
特に、$\operatorname{Gr}(4,9)$ 上の階数十八束の零点という形式と随伴公式だけから、
$X_Q$ 全体を滑らかなK3曲面とみなすことはできない。
閉包の全既約成分、交差スキーム、正規交差性はこの議論では決定していない。

\paragraph{補助スクロールの次数.}
$E_Q$ が滑らかな場合、$\mathcal R=\mathcal O_{E_Q}^6/\mathcal K_Q$ は階数四で、
$\mathcal O_{E_Q}(1)$ に値を持つ非退化な交代形式を持つ。
従って $\det\mathcal R=\mathcal O_{E_Q}(2)$、
$\det\mathcal K_Q=\mathcal O_{E_Q}(-2)$ である。
$\mathbb P(\mathcal K_Q)$ を直線部分空間をパラメータ化する規約で取ると、
自然な超平面類 $H=c_1(\mathcal O_{\mathbb P(\mathcal K_Q)}(1))$ に対して
\[
 H^2=-\deg\mathcal K_Q=6.
\]
一般の二つの相補的な $U_\pm$ を用いれば、核束はその二つへの交わりの線束の直和に分かれる。
これらは、境界平面と楕円スクロールの閉包との接合を調べる際の有用な出発点になる。
